\section{总结与展望}
\label{sec:conclusion}

\subsection{工作总结}

本文针对空地协同多功能雷达/电子对抗网络的部署优化问题，基于MOPSO-DT算法框架，系统地研究了复杂区域内的多目标部署优化方法。主要工作总结如下：

\textbf{（1）边界效应的系统抑制。}将复杂区域凸分解与坐标变换相结合，把受物理边界约束的优化问题转化为归一化空间中的无约束搜索。该方法避免了外接矩形采样后剔除非法点带来的可行域欠采样问题，使搜索过程能够覆盖部署区域的完整可行空间。

\textbf{（2）传播特性的差异化建模。}引入区分空-地（A2G）和地-地（G2G）链路的异构传播模型，并兼容雷达方程精细模型与指数衰减简化模型两种层级，适应从工程分析到快速原型的不同应用需求。

\textbf{（3）算法参数的系统优化。}通过6个维度的调优实验确定了推荐配置（standard + crowding + $p_m=0.01$），超体积（HV）较原始配置提升5.9\%，Pareto解分布均匀性（Spacing）改善54\%，Numba JIT加速使完整优化（$T_{\max}=500$）耗时仅211.9秒。

\textbf{（4）ECR-$J_{\min}$权衡关系的定量揭示。}在两种不同的传播模型和场景规模下，实验均观察到有效覆盖率（ECR）与最小等效干扰强度（$J_{\min}$）之间存在强负相关关系（$r < -0.93$，最高$r = -0.958$）。这一结果表明，在有限节点数量和距离衰减约束下，广域覆盖与最薄弱任务点压制强度之间存在显著经验性权衡，可为部署方案选择提供定量参考。

\subsection{研究局限性}

本文工作存在以下局限性，需要在后续研究中加以改进：

\begin{itemize}
    \item \textbf{仿真环境理想化。}所有实验基于二维平面模型，未考虑真实的三维地形遮挡。实际战场环境中，高程变化、建筑物遮挡等因素将显著影响电磁波传播。

    \item \textbf{传播模型简化。}当前传播模型采用距离依赖的指数衰减或自由空间雷达方程，未引入多径效应、阴影衰落等随机因素。在复杂城区或山地环境中，确定性的传播模型可能产生较大偏差。

    \item \textbf{单次运行统计。}受计算资源限制，当前实验均为单seed运行。粒子群优化作为随机搜索算法，其结果具有一定随机性。多seed重复实验和统计显著性检验将增强结论的可靠性。

    \item \textbf{权衡结论的外推范围有限。}ECR与$J_{\min}$的强负相关关系来自本文两个主要实验场景，能够支持当前场景下的部署权衡分析，但尚不能证明其在所有雷达数量、任务分布、传播参数和约束条件下均成立。

    \item \textbf{静态部署假设。}本文假设雷达位置在部署后固定不变，未考虑目标移动性对部署方案的影响。在动态战场环境中，最优部署方案可能需要随敌我态势变化而实时调整。

    \item \textbf{基准对比有限。}NSGA-II和MOEA/D等经典多目标算法尚未在相同场景下直接对比。更完整的算法对比将有助于全面评估MOPSO-DT的优劣势。
\end{itemize}

\subsection{未来研究方向}

基于上述局限性，未来工作可从以下方向展开：

\begin{itemize}
    \item \textbf{动态场景扩展。}将静态部署模型扩展为动态重部署问题，考虑目标移动轨迹预测和在线优化决策，研究基于滚动时域优化的实时部署调整策略。

    \item \textbf{精细传播模型。}引入数字高程模型（DEM）和射线追踪方法，实现对地形遮挡和非视距传播的精确建模。结合实测信道数据校准模型参数。

    \item \textbf{多seed统计分析。}对每个配置进行至少5次独立重复实验，报告均值和标准差或95\%置信区间，提升结论的统计严谨性。

    \item \textbf{算法对比扩展。}在相同场景和评估协议下，实现NSGA-II、MOEA/D、SPEA2等经典多目标算法，进行系统性的公平对比。

    \item \textbf{半实物仿真验证。}结合软件无线电（SDR）平台和信道模拟器，在受控环境中验证算法生成的部署方案在实际电磁环境中的表现。

    \item \textbf{三维部署扩展。}将二维平面部署扩展至三维空间，考虑空中节点的飞行高度优化，进一步提升空地协同网络的覆盖和压制效能。
\end{itemize}
